\documentclass[11pt]{article}
\usepackage[T1]{fontenc}
\usepackage[utf8]{inputenc}
\usepackage{amsmath,amssymb,amsthm,booktabs}
\usepackage[margin=1in]{geometry}
\usepackage{hyperref,xurl}
\providecommand{\tightlist}{\setlength{\itemsep}{0pt}\setlength{\parskip}{0pt}}
\title{The equality structure of the second-volume\linebreak ideal right-angled hyperbolic four-polytope}
\author{Coxeter Polytopes Frontier Project}
\date{Proof translation, 18 October 2026}
\begin{document}
\maketitle
\begin{abstract}
We classify the equality case in the internally proved second-volume theorem for entirely ideal right-angled hyperbolic four-polytopes. Every such polytope of volume $8\pi^2/3$ is the union of two ideal right-angled 24-cells reflected across a common octahedral facet. The proof first forces eight doubled-octahedron facets. Their six nonsquare cusps determine an octahedral incidence graph, and equal horospherical widths force a common cutting hyperplane. Both halves then have only octahedral facets and are 24-cells. Published three-dimensional facet classifications and rigidity are explicit inputs; the new incidence, normal and cutting arguments are internal. Exact finite certificates support the model calculations but do not replace the universal proof. Three independent agent reviews have passed; external refereeing and a complete Lean formalization are not claimed.
\end{abstract}
18 October 2026. \textbf{Internally checked universal equality theorem, with three independent full written-proof reviews and exact supporting certificates.} This is separate from the frozen \href{ideal_ra4_second_volume_20261018.md}{second-value theorem}. The earlier theorem and construction certificate remain unchanged. External human review and end-to-end Lean formalization are not claimed.

\textbf{Theorem.} Every finite-volume entirely ideal right-angled hyperbolic four-polytope of volume \(8\pi^2/3\) is isometric to the double of the ideal right-angled 24-cell across an octahedral facet. In particular its f-vector is \((42,168,164,38)\), and its facets are thirty regular ideal octahedra and eight doubles of such octahedra.

The argument below is universal. A small supporting incidence/normal certificate can check its finite model facts, but neither a bounded four-polytope census nor equality in the earlier coarse lower bound is assumed. In fact the known double does not make all those coarse inequalities equalities.

\section{Inputs and notation}\label{inputs-and-notation}

Use the frozen proof's identities

\[
 W=3(v-f),\qquad \mathrm{Vol}=4\pi^2/3+W\pi^2/9,
 \qquad d(F)=v_F-6=\sum_{R\subset F}(k(R)-3),
 \qquad \sum_Fd(F)=2W.
\]

At this volume, \(W=12\) and \(\sum_Fd(F)=24\). Call a facet \textbf{positive} if its defect is positive. All other facets are regular ideal octahedra. Every positive defect is at least 2. Write \(N\) for the number of positive facets and \(m\) for their minimum defect. For any positive facet \(F\), write q for its nontriangular faces and s for its nonsquare cusp rectangles. The already proved opposite-cusp injection gives the useful non-minimal form

\[
 24\ge d(F)+m(q+s). \tag{1}
\]

The published antiprism bound and equality case, and the small admissible facet list, are taken from \href{https://msp.org/agt/2012/12-4/agt-v12-n4-p03-s.pdf}{Kolpakov (2012)}, Proposition 5 and §4.3, pp.1951--1952. In addition to the nine-vertex entry used in the gap proof, this equality argument uses the ten-vertex \textbf{octahedrite} entry in Figure 7. Thus this is an explicit additional published input. It is the unique admissible ten-vertex facet having only triangles and quadrilaterals. Its four quadrilateral faces have edge-adjacency graph consisting of two disjoint edges. The ten-vertex pentagonal antiprism is handled separately.

We use right-angled rigidity in dimension three to identify these classified facets geometrically. Denote the unique nine-vertex facet by D9. It is a double octahedron. Its three quadrilateral faces have no common edges, and each pair meets at one of the three nonsquare cusps. Those three cusps form the internal gluing triangle. Each quadrilateral contains exactly two of them, as opposite vertices; their connecting segment is the interior gluing diagonal. Every triangular face contains at most one nonsquare cusp. At a nonsquare cusp the rectangle has side ratio two, and \textbf{the two quadrilateral sides are its two long sides}: the rectangle is obtained by joining two squares, and the merged face edges are precisely the long sides. At each of the other six, square cusps exactly one of the four incident polygonal faces is a quadrilateral. These facts follow by unfolding two octahedra and do not require a four-dimensional classification.

\section{A basic acute-intersection fact}\label{a-basic-acute-intersection-fact}

We use this internal fact to pass from pairwise facet adjacency to actual common faces. Give every supporting wall a unit outward spacelike normal. Distinct wall products are nonpositive for an acute Coxeter chamber. If \(n_1,\dots,n_k\) are mutually orthogonal and p is a strict future timelike interior point, put

\[
 u=p-\sum_{i=1}^k\langle p,n_i\rangle n_i.
\]

Then \(u^2=p^2-\sum_i\langle p,n_i\rangle^2<0\), \(\langle u,n_i\rangle=0\), and every other wall product is strictly negative: each subtracted product \(\langle p,n_i\rangle\langle n_i,n_j\rangle\) is nonnegative. The projected timelike vector has the future sign, for example because \(\langle p,u\rangle=u^2<0\). Thus the k walls meet in an actual codimension-k face. Three mutually adjacent right-angled facets meet along an actual edge; four would give a finite vertex, forbidden here.

Consequently, if F has nontriangular neighbors G and H and G is adjacent to H anywhere in P, the corresponding two faces of F share an edge. This excludes an adjacency elsewhere that is invisible in the intrinsic facet picture.

\section{First reduction: the minimum positive defect is three}\label{first-reduction-the-minimum-positive-defect-is-three}

\textbf{No antiprism \(A_k\) with \(k\ge5\).} Its defect is \(d=2k-6\ge4\), \(q=2\) and \(s=2k=d+6\). Even the universal \(m\ge2\) in (1) gives

\[
 24\ge d+2(d+8)=3d+16\ge28,
\]

a contradiction.

\textbf{No defect two facet.} Such a facet is the square antiprism A4. It has eight nonsquare cusps and two quadrilateral faces; each cusp belongs to exactly one of these quadrilaterals. If one face of a rectangular3-box is nonsquare, at least one of the other two opposite-face pairs is nonsquare. Otherwise both ratios to the third side would be one and all sides would be equal. Hence at each cusp of A4 at least two of its four ridge-neighbor facets are nonoctahedral. At most one is the unique quadrilateral neighbor through that cusp. At least one triangular neighbor is therefore positive. Covering the eight cusps by triangular faces requires at least three distinct such triangular neighbors. These are distinct from F, its two quadrilateral neighbors, and its eight nonsquare-cusp opposites, by the opposite-cusp injection. Thus at least 14 positive facets exist, giving total defect at least 28. This is impossible. We have \(m\ge3\).

\textbf{The branch \(m\ge5\) is impossible.} There are at most four positive facets. Every positive facet is a non-antiprism, so has at least three nontriangular faces. Their neighbors must be distinct positive facets. Thus there are exactly four, all pairwise adjacent. The acute-intersection fact would give a finite vertex.

\textbf{The branch \(m=4\) is impossible.} There are at most six positive facets. Choose a defect four facet \(F\). A pentagon would force A5, already excluded. Thus F is the Figure 7 ten-vertex octahedrite, with four quadrilaterals paired by edge adjacency. Let \(G_1,\dots,G_4\) be their four positive neighbors. The induced adjacency graph on these four neighbors is contained in that matching, by the acute-intersection fact. Each \(G_i\) has at most F, its one matching partner and one further positive facet as possible nontriangular neighbors. Every \(G_i\) needs at least three such neighbors, so a further facet must exist and \(q(G_i)=3\). None of the \(G_i\) can have defect four, since that case has four nontriangular faces; hence each has defect at least 5. Including F and the necessary further facet already gives \(4+4\cdot5+4=28\), a contradiction.

It follows that \textbf{\(m=3\)}, so some facet is D9.

\section{Exactly eight positive facets, all D9}\label{exactly-eight-positive-facets-all-d9}

A D9 facet and its three quadrilateral neighbors and three nonsquare-cusp opposites give seven positive facets. Since \(m=3\) and the total defect is 24,

\[
 7\le N\le8. \tag{2}
\]

If \(N=8\) all defects are3, as desired. Suppose \(N=7\). For every D9 facet, the seven just counted exhaust all positive facets. In particular:

\begin{quote}
Under \(N=7\), no triangular ridge of a D9 facet has a positive neighbor.
\end{quote}

The seven defects are at least3 and sum24. The possible excess partitions above seven copies of3 are: one defect six; a defect5 and a defect four; or three defects4. There are at most three facets whose defects exceed3.

There can be no defect four facet \(F\). Its four quadrilaterals occur in edge-adjacent pairs. Suppose one quadrilateral neighbor G were D9, and let H be the neighbor across its paired quadrilateral. The edge shared by the two quadrilaterals of F is an actual edge shared by F,G,H. In G, the faces cut out by F and H share this edge. The former is a quadrilateral. The latter cannot be a quadrilateral because D9's quadrilaterals share no edges, and cannot be a triangle because H is positive and \(N=7\) forbids a positive triangular neighbor. Thus every one of F's four quadrilateral neighbors would have defect greater than3. There are at most two such facets other than F, a contradiction.

We are left with one defect six facet \(F\) and six D9 facets. Every nontriangle of F is shared with a D9, and is therefore a quadrilateral. There are exactly six quadrilaterals. No two share an edge, by the same argument with their common D9 neighbor: two adjacent quadrilaterals of F would force either adjacent quadrilaterals or a positive triangular neighbor in D9.

The facet \(F\) has twelve vertices. Its six quadrilaterals give24 quadrilateral-vertex incidences. At a four-valent vertex no two consecutive faces can both be quadrilaterals; therefore there are at most two quadrilaterals there. Equality in the incidence count forces exactly two, alternating with two triangles, at every vertex. Form a graph J whose vertices are the six quadrilateral faces and whose edges are the twelve vertices of F, joining the two incident quadrilaterals. It is a simple 4-regular graph on six vertices: two distinct quadrilateral faces cannot share two vertices without sharing an edge, since their intersection is a face of the convex3-polyhedron. Its complement is a perfect matching, so it is the octahedral graph.

The embedding and medial recovery are explicit here. Draw each new edge from a center of one quadrilateral to their shared original vertex and then to the center of the other quadrilateral. Alternation at every original vertex makes this a spherical embedding, with the eight original triangular regions as its eight triangular faces. Those face cycles are distinct: a connected embedded graph with more than three vertices cannot have both sides of one3-cycle as faces, since then no other vertex or edge could lie on either side. The octahedral graph has exactly eight 3-cycles, so all its 3-cycles are precisely these faces. This determines the face incidences, not merely the abstract graph. Recover F with one vertex for each edge of J, its six quadrilaterals around the six vertices of J, and its eight triangles around those eight 3-cycles. This is exactly the cuboctahedral face lattice. No general embedding-rigidity theorem or unbounded facet census is being assumed.

The ideal right-angled cuboctahedron has all twelve cusp rectangles nonsquare, of aspect \(\sqrt2\). An exact model in Lorentz coordinates has vertices

\[
 (1,x),\quad x\text{ a permutation of }(\pm1,\pm1,0)/\sqrt2;
\]

its six square-wall normals are \((1,\pm\sqrt2 e_i)\), of norm one, and its eight triangle-wall normals are \((1,(\pm1,\pm1,\pm1)/\sqrt2)\), of norm \(1/2\). At \((1,1/\sqrt2,1/\sqrt2,0)\), the two opposite square normals sum to \((2,\sqrt2,\sqrt2,0)\), and the two opposite \emph{unnormalized} triangle normals have the same sum. Unit normalization multiplies the latter by \(\sqrt2\). The horospherical width ratio is therefore \(\sqrt2\), and symmetry gives it at every vertex. The model is the cube \(|X_i|\le1\) truncated by \(\sum|X_i|\le2\) in scaled Klein coordinates, so its combinatorics are explicit. Right-angled rigidity identifies F with this model.

The opposite-cusp injection for F would now give \(1+6+12=19\) positive facets, contradicting \(N=7\). We conclude:

\[
 \boxed{N=8,\quad d(F)=3\text{ for every positive facet}.} \tag{3}
\]

All other facets are regular ideal octahedra. All nontriangular ridges are quadrilaterals; there are twelve of them.

\section{The six nonsquare cusps and their exact incidence graph}\label{the-six-nonsquare-cusps-and-their-exact-incidence-graph}

Each D9 contributes three nonsquare cusp rectangles, for 24 incidences. At an ambient cusp every face rectangle belongs to an octahedron or D9, and hence has aspect one or two. If its box has three distinct side lengths, their pairwise ratios cannot all lie in \(\{1,2\}\). Thus every noncubic cusp box has side ratios either \(2:1:1\) or \(2:2:1\).

The second possibility is impossible. In a \(2:2:1\) box the four nonsquare faces must be D9 faces. In each such face, its two long sides are quadrilateral ridges. Consequently each square \(2\times2\) end face has four quadrilateral sides incident at the cusp. It cannot be an octahedron, which has none, or a D9 at a square cusp, which has exactly one. Thus every noncubic cusp box has side ratio \(2:1:1\).

For this box the four nonsquare D9 faces form a belt around the long direction, and their four belt adjacencies are quadrilateral ridges. The two square end faces meet the D9 belt along short, triangular sides. They cannot themselves be D9 square cusps, because all four of their incident sides are triangular rather than one quadrilateral; hence they are octahedral. Each noncubic cusp therefore accounts for exactly four of the24 nonsquare D9 incidences. There are exactly six such cusps.

Define a graph J with these six cusps as vertices. A quadrilateral ridge supplies an edge joining its two nonsquare cusps. They are distinct, as seen intrinsically in D9. Every cusp lies on exactly four quadrilateral ridges, so J is4-regular. It is \textbf{simple}: if two distinct cusps belonged to two distinct quadrilateral ridges, the union of the incident D9 facets of those ridges would contain at least three D9 facets. Among three of the four D9 belt facets at a cusp, two are opposite box faces. Their supporting hyperplanes are asymptotic at that one cusp and cannot contain a second ideal point. This contradiction excludes parallel edges. (Two distinct ridges cannot have the same pair of incident facets.)

Thus J is a simple 4-regular graph on six vertices: it is \(K_6\) minus a perfect matching, the octahedral graph. This deduction does not assume that an arbitrary cubic eight-vertex graph is the cube, and uses no unproved surface or planarity claim.

The three nonsquare cusps of each D9 form a triangle in J, because its three quadrilaterals connect each pair. Distinct D9 facets give distinct triangles: otherwise they both contain the convex hull of the same three ideal points, but this triangle is the internal gluing triangle in the relative interior of either D9, so it cannot lie in another supporting facet. There are eight D9 facets and exactly eight triangles in J. Hence all eight triangles occur. Label the three nonedges of J as three opposite vertex pairs. A D9 is then labeled by a sign vector \(s\in\{\pm1\}^3\), choosing one cusp from each pair. Two D9 facets share a quadrilateral precisely when their signs differ in one coordinate. Their quadrilateral adjacency graph is therefore the cube.

\section{Recovering the common gluing hyperplane}\label{recovering-the-common-gluing-hyperplane}

Choose outward \textbf{unit} normals \(n_s\) for the eight D9 facets. At a \(2:1:1\) cusp, the four belt normals form two opposite pairs, and both pairs bound short directions of equal width. For opposite unit normals their sum is a positive multiple of the cusp null vector, with coefficient proportional to the corresponding width. Therefore the sums of the two opposite pairs are equal. In the sign labeling, these are precisely the parallelogram equations on all six square faces of the cube.

The cube equations imply

\[
 n_s=a+s_1b_1+s_2b_2+s_3b_3. \tag{4}
\]

For example, take a and the b's to be the constant and linear Fourier averages of the eight n's; every mixed finite difference vanishes by a square equation, so there are no higher terms. All eight unit-norm equations give \(\langle a,b_i\rangle=0\), \(\langle b_i,b_j\rangle=0\) for \(i\ne j\), and \(a^2+\sum b_i^2=1\). The right-angle equation for each cube edge gives \(1-2b_i^2=0\). Thus

\[
 a^2=-\tfrac12,\qquad b_i^2=\tfrac12,
 \qquad a,b_1,b_2,b_3\text{ mutually orthogonal}. \tag{5}
\]

They span a Lorentz 4-space U. Its orthogonal complement in the ambient Lorentz 5-space is one spacelike line. Let H be the hyperbolic 3-space corresponding to U. All eight D9 wall hyperplanes are orthogonal to H.

In orthonormal coordinates in U, the eight normals are \((1,s_1,s_2,s_3)/\sqrt2\). The time sign is future, as follows also from the positive cusp-null sums. Their halfspaces in H enclose the regular ideal octahedron

\[
 S=\{(t,x)\in\mathbb H^3:|x_1|+|x_2|+|x_3|\le t\}.
\]

The six ideal vertices of S are the six actual nonsquare cusps: each opposite belt-pair sum is \(2(a\pm b_i)\). Since P contains these six cusps, convexity gives \(S\subset P\). Conversely the eight D9 inequalities give \(P\cap H\subset S\). Hence

\[
 P\cap H=S. \tag{6}
\]

The face of S in any D9 is the convex hull of its three nonsquare cusps, exactly that facet's internal gluing triangle. Thus H cuts every D9 into its two regular ideal octahedra. The relative interior of every face of S is in the relative interior of a D9 facet. The relative interior of every edge of S is the internal diagonal of an actual quadrilateral ridge, not an edge of P. This follows from the intrinsic D9 model: the edge joins the two nonsquare vertices of the quadrilateral, which are opposite vertices after the two triangles are joined.

\section{The section cuts P into two 24-cells}\label{the-section-cuts-p-into-two-24-cells}

Let O be any octahedral facet of P. It cannot meet S at a finite point. Indeed a zero set of its supporting functional on S is a face of S (or all S). If it contained a finite point in the relative interior of S, it would contain S; if it met a face interior, it would contain that face. Both are impossible because the interiors of the triangular faces of S lie inside other facets D9. If it met the interior of an edge of S, it would contain a relative-interior point of a quadrilateral ridge, whose only incident facets are its two D9 facets. Thus \(O\cap S\) contains at most ideal vertices, and in fact at most one, since a segment between two would contain finite points.

If O had points on both sides of H, their geodesic segment in O would cross H at a finite point of \(P\cap H=S\), contradicting the preceding paragraph. Therefore no octahedral facet is cut by H. Only the eight D9 facets are cut, each along its internal triangle. Both sides are nonempty, since those triangles are internal to the D9 facets.

The two intersections \(P_\pm=P\cap H^\pm\) are finite-volume entirely ideal right-angled four-polytopes. Their facets are the uncut original octahedra on the respective side, the octahedral halves of the D9 facets, and the new octahedral facet S. They acquire no finite vertices: the cut section has only the six already ideal vertices, and the explicit facet cuts introduce no other vertices. New finite dihedral intersections with S involve only D9 walls, already orthogonal to H. All other dihedral angles are inherited.

Both halves have only triangular ridges, hence W0. The published unique-minimum theorem makes each half an ideal right-angled 24-cell. They share S and occupy opposite sides. Here is the precise placement argument. A regular 24-cell is flag-transitive. Any symmetry of its octahedral facet sends a flag in that facet to another; flag-transitivity of the cell supplies a cell symmetry with that same flag action and stabilizing the facet. An isometry of the regular octahedron fixing a complete flag is the identity, so this restriction is the prescribed facet symmetry. Thus the facet stabilizer surjects onto the full octahedral symmetry group. Given an isometry from one half to the reflection of the other, first align the distinguished facets by cell symmetry and then cancel its restriction to S by such a stabilizer. The resulting isometry fixes S pointwise. An ambient hyperbolic isometry fixing a hyperplane pointwise is either the identity or reflection in it; the same-side condition selects the identity. Therefore the halves are precisely reflections across S, up to ambient isometry.

This proves the equality theorem. The f-vector and exact facet/cusp counts follow from the frozen double certificate. No equality theorem is inferred solely from the eight-facet defect count, and no graph realization or geometric gluing step has been omitted into a computational census.

\section{Scope and review ledger}\label{scope-and-review-ledger}

The new published input is the complete ten-vertex octahedrite entry, with its four-quad matching incidence read from Figure 7. The cuboctahedron reduction, nonsquare box propagation, simple six-cusp graph, unit-normal cube equations and cutting argument are internal proof steps requiring review. The old antiprism shape and opposite-cusp injection are inherited from the frozen gap proof, not relabeled as external.

Root's independent full written-proof mathematical review passed, including the explicit embedded medial recovery and flag-transitive final placement. It independently checked the Figure 7 quadrilateral matching directly against the primary PDF image and its explicit complete-list statement. Its full supporting producer and standalone-helper source reads also passed. It performed no certificate execution or historical classification reproof.

The supporting producer \href{../scripts/fresh_ideal_ra4_equality_structure_20261018.py}{fresh\_ideal\_ra4\_equality\_structure\_20261018.py} initially generated its exact certificate in run \texttt{373801}; separate replay \texttt{ad43d5} passed. It checks the Figure 7 incidence transcription, D9 local face/cusp data, the complete cuboctahedral 364-active-triple domain and all 12 cusp-width squares, all 15 labeled simple 4-regular six-vertex graphs, and the cube normal algebra. The \href{../scripts/check_ideal_ra4_equality_structure_saved_20261018.py}{standalone checker}, initial run \texttt{83f326}, independently checks these finite supports with reverse-column elimination and imports no producer. Final generation, a separate-process fresh replay, and the hardened helper all passed in terminal run \texttt{bb139d}, after adding the producer's refusal to overwrite an existing certificate and the helper's explicit cuboctahedral-wall-set guard. The previous unfrozen JSON was preserved under \texttt{/tmp} before that source-bound regeneration. Such a certificate does not substitute for the universal reductions and convexity arguments above.

Order16's independent full written-proof review passed, including the embedded medial recovery and flag-transitive final placement. It also reread the frozen gap proof, the published completeness text and the primary Figure 7 image, confirming the quadrilateral matching. Its scope includes every minimum-defect and \(N=7\)/\(N=8\) case, the box-ratio restriction, simple six-cusp graph, cube-normal algebra, actual section and octahedral no-crossing argument. It did not execute the certificate, reprove the historical classification, or read the new supporting producer/helper; those scopes are not attributed to it.

Construction's third independent full written-proof mathematical read passed, including a fresh primary-text and visual Figure 7 check. It checked all universal reductions, the cuboctahedron model, cusp-box propagation, six-cusp graph, affine normal equations and exact cutting argument. It did not execute geometry or fully review the new supporting producer/helper. Its formatting correction was applied: all mathematical delimiters and angle-bracket escapes are now ordinary valid Markdown/TeX, with no embedded control characters. The prior second-value report remains immutable; a separate equality manuscript translates this new theorem.

\end{document}
